\documentclass{article}
\usepackage{graphicx}
\usepackage{amsmath}
\usepackage{amsfonts}
\usepackage{amssymb}

\title{Optimal Trade-Off-Query Algorithmus}
\author{Jan Oberwahrenbrock}
\date{April 2026}

\begin{document}

\maketitle

\section*{Input}

\begin{itemize}
\item
Handlungsalternativenmatrix
\[
H=(u_{k\ell}) \in [0,1]^{m\times n},
\]
wobei \(u_{k\ell}\) den Nutzenwert der Alternative \(k\) im Ziel \(\ell\)
bezeichnet.

\item
Menge der bisher beantworteten Trade-off-Queries
\[
T=\{(a,b,\xi,\circ)\},
\]
wobei \(a,b\) Zielindizes, \(\xi\ge 0\) ein Schwellenwert und
\[
\circ \in \{<,=,>\}
\]
der zugeh\"orige Operator ist. Eine solche Query hat die Form
\[
w_a \ \circ \ \xi\, w_b.
\]
\end{itemize}


\section*{Ablauf}

\begin{enumerate}
\item
Bestimme aus \(T\) die aktuelle Zielgewichtsmenge \(W(T)\).

\item
Bestimme f\"ur jede Alternative \(k\) die Optimalit\"atsregion
\[
O_k(T)
=
\left\{
w\in W(T)\;\middle|\;
G_k(w)\ge G_g(w)\ \text{f\"ur alle Alternativen } g
\right\},
\]
wobei \(G_k(w)=\sum_{\ell=1}^n u_{k\ell}w_\ell\) den Gesamtnutzen der Alternative \(k\) bezeichnet.


Bestimme daraus die Kandidatenmenge
\[
K(T)=\left\{k\mid O_k(T)\neq\emptyset\right\}.
\]

\item
Pr\"ufe, ob einer der folgenden Terminierungsf\"alle eintritt:
\[
|K(T)|=1
\]
oder
\[
G_a(w)=G_b(w)\qquad\text{f\"ur alle }a,b\in K(T)\text{ und alle }w\in W(T).
\]

Im zweiten Fall kann keine weitere Query die verbleibenden Kandidaten noch
unterscheiden.

\item
Bestimme f\"ur jedes geordnete Zielpaar \((i,j)\) mit \(i\neq j\) und jeden
Kandidaten \(k\) das Verh\"altnisintervall
\[
I_k^{(i,j)}
=
\left\{
\frac{w_i}{w_j}
\;\middle|\;
w\in O_k(T),\; w_j>0
\right\}
=
[\alpha_k^{(i,j)},\beta_k^{(i,j)}].
\]

Bestimme daraus f\"ur jedes geordnete Zielpaar \((i,j)\) die Menge der linken
Rohkandidaten
\[
R_{\mathrm{left}}^{(i,j)}
=
\left\{
(i,j,\alpha_k^{(i,j)},\mathrm{left})
\;\middle|\;
k\in K(T),\ \alpha_k^{(i,j)} > \varepsilon
\right\}
\]
sowie die Menge der rechten Rohkandidaten
\[
R_{\mathrm{right}}^{(i,j)}
=
\left\{
(i,j,\beta_k^{(i,j)},\mathrm{right})
\;\middle|\;
k\in K(T),\ \beta_k^{(i,j)} < \infty
\right\}.
\]

Definiere daraus f\"ur jedes geordnete Zielpaar \((i,j)\) die
Rohkandidatenmenge
\[
R^{(i,j)}
=
R_{\mathrm{left}}^{(i,j)}
\cup
R_{\mathrm{right}}^{(i,j)}
\]
sowie die gesamte Rohkandidatenmenge
\[
R
=
\bigcup_{i\neq j} R^{(i,j)}.
\]

Fasse anschlie\ss end \"aquivalente Rohkandidaten zusammen, wobei f\"ur \(x>0\)
\[
(i,j,x,\mathrm{left}) \sim (j,i,1/x,\mathrm{right})
\]
und
\[
(i,j,x,\mathrm{right}) \sim (j,i,1/x,\mathrm{left})
\]
gelten.



Die daraus entstehende gefilterte Rohkandidatenmenge bezeichne mit
\[
R^{\mathrm{filtered}}.
\]

Definiere daraus die Mengen
\[
Q_{\mathrm{left}}^{\mathrm{filtered}}
=
\left\{
(i,j,\alpha-\varepsilon)
\;\middle|\;
(i,j,\alpha,\mathrm{left})\in R^{\mathrm{filtered}}
\right\}
\]
und
\[
Q_{\mathrm{right}}^{\mathrm{filtered}}
=
\left\{
(i,j,\beta+\varepsilon)
\;\middle|\;
(i,j,\beta,\mathrm{right})\in R^{\mathrm{filtered}}
\right\}.
\]

Die gefilterte Query-Menge ist dann gegeben durch
\[
Q^{\mathrm{filtered}}
=
Q_{\mathrm{left}}^{\mathrm{filtered}}
\cup
Q_{\mathrm{right}}^{\mathrm{filtered}}.
\]





\item
Bestimme f\"ur jede Query \(q=(i,j,x)\in Q^{\mathrm{filtered}}\) die Mengen
\[
K_{<}(q):=\{k\in K(T)\mid x>\alpha_k^{(i,j)}\},
\]
\[
K_{>}(q):=\{k\in K(T)\mid x<\beta_k^{(i,j)}\}.
\]

Bestimme zus\"atzlich die Antwortwahrscheinlichkeiten
\[
p_{<}(q),\qquad p_{>}(q),
\]
also die Wahrscheinlichkeiten daf\"ur, dass auf die Query \(q\) eine
\(<\)- bzw. eine \(>\)-Antwort erfolgt.

\item
Definiere die Menge der informativen Queries durch
\[
Q^{\mathrm{info}}
=
\left\{
q\in Q^{\mathrm{filtered}}
\;\middle|\;
\bigl(K_{<}(q)\subsetneq K(T)\ \text{und}\ p_{<}(q)>0\bigr)
\ \lor\
\bigl(K_{>}(q)\subsetneq K(T)\ \text{und}\ p_{>}(q)>0\bigr)
\right\}.
\]


Falls
\[
Q^{\mathrm{info}}=\emptyset,
\]
terminiert der Algorithmus, da keine weitere informative Query existiert.

\item
Falls \(Q^{\mathrm{info}}\neq\emptyset\), bestimme f\"ur jede Query
\(q\in Q^{\mathrm{info}}\) die erwartete verbleibende Kandidatenanzahl
\[
\mathbb{E}[N\mid q]
=
p_{<}(q)\,|K_{<}(q)|
+
p_{>}(q)\,|K_{>}(q)|.
\]

W\"ahle anschlie\ss end eine Query
\[
q^\star \in \arg\min_{q\in Q^{\mathrm{info}}} \mathbb{E}[N\mid q].
\]

\end{enumerate}

\end{document}
